\begin{theorem}[I8: the full presentation in the finite-topology category; Proved]\label{i:linearuniversal}
Let $A$ be a regular noetherian local ring of finite dimension, and
let $K$ be its residual resolution, normalized by height. Put
$B=\prod_{\mathfrak p}\widehat{A_{\mathfrak p}}$.
The finite-topology algebra $\End_A^*(K)$ is universal among
separated complete graded topological $A$-algebras whose topology
has a basis of open left ideals, equipped with the prescribed
continuous degree-zero map from $B$ and an element $T$ of degree one
satisfying
\[
 T^2=0,\qquad e_{\mathfrak q}Te_{\mathfrak p}=0
 \quad\text{unless }\mathfrak p\subset\mathfrak q
 \text{ is an immediate specialization}.
\]
In particular the completed presentation in this category is an
isomorphism in every degree and in every dimension. The universal
property allows a continuous map from $B$; requiring it to identify
degree zero gives the version in the question.
\end{theorem}

\begin{proof}
We first prove a statement about discrete representations that does
not assume the presentation.

\emph{Cousin complexes.}
Call a complex $M$ a Cousin complex here if
\[
 M^i=\bigoplus_{\operatorname{ht}\mathfrak p=i}M_{\mathfrak p},
\]
where $M_{\mathfrak p}$ is an $A_{\mathfrak p}$-module all of whose
elements are killed by powers of $\mathfrak pA_{\mathfrak p}$.
Its differentials are $A$-linear. A map between two such point
modules is zero unless the source prime is contained in the target
prime, by the element-annihilator argument of
Proposition~\ref{i:zero}. Thus the differential has the required
incidence property. The action of
$\widehat{A_{\mathfrak p}}$ on $M_{\mathfrak p}$ is canonical:
on each element it is determined modulo a power of
$\mathfrak pA_{\mathfrak p}$. Every $A$-linear morphism of Cousin
complexes commutes with these completed diagonal actions.

\emph{A universal property of a complete-intersection Cousin
resolution.}
Let $J=(x_1,\ldots,x_h)$ be generated by an $A$-regular sequence,
put $N=A/J$, and put
\[
 C_J=\Hom_A(N,K)=K[J].
\]
The Koszul resolution and its self-duality give
$\Ext_A^i(N,A)=0$ for $i\ne h$ and
$\Ext_A^h(N,A)\simeq N$
\cite[Lemma 48.15.5, tag 0BQZ]{StacksKoszulDuality}.
Choose this last identification, so that
$N$ identifies with $Z^h(C_J)$. Every prime containing $J$ has
height at least $h$, so $C_J^i=0$ for $i<h$. In degree $h$ its
terms identify with
\[
 C_J^h=\bigoplus_{\substack{\mathfrak r\supset J\\
                            \operatorname{ht}\mathfrak r=h}}
 N_{\mathfrak r},
\]
with the augmentation $N\to C_J^h$ given by localization.
Indeed $N$ is Cohen--Macaulay, has no embedded primes, and the
Koszul self-duality identifies the localized top Ext at each
minimal prime with the corresponding Artinian complete intersection.

For every Cousin complex $M$, an $A$-linear map
$N\to Z^h(M)$ extends uniquely to a morphism
\[
                         C_J\longrightarrow M.          \tag{I8.1}
\]
Here is the full extension argument. In degree $h$, localize the
given map at each height-$h$ prime and take their direct sum. Suppose
degrees through $j-1$ have been defined compatibly with the
differentials, where $j>h$. For a prime $\mathfrak q$ of height
$j$, project the composite
$C_J^{j-1}\to M^{j-1}\to M^j$ to $M_{\mathfrak q}$ and
localize at $\mathfrak q$. It kills the kernel of
\[
 (C_J^{j-1})_{\mathfrak q}\longrightarrow
 (C_J^j)_{\mathfrak q}=C_{J,\mathfrak q}^j.
\]
For $j=h+1$ this follows because the prescribed map on $N$ is a
cycle; for larger $j$ it follows from the already established chain
identity and $d_M^2=0$. The displayed arrow is surjective because
the localized complex has no cohomology in degree $j>h$ and no
terms above $j$. Consequently the composite factors uniquely
through $C_{J,\mathfrak q}^j$. Take the direct sum over
$\mathfrak q$. Each source element has finitely many components,
so this gives a well-defined morphism, and comparison of all target
components proves the chain identity. Surjectivity in this
construction also proves uniqueness at every stage. This proves
(I8.1).

\emph{These resolutions generate every Cousin complex.}
Take $v\in M_{\mathfrak p}$ and set
$h=\operatorname{ht}\mathfrak p$. Its differential has finitely
many nonzero components, at primes $\mathfrak q$ strictly containing
$\mathfrak p$. For each such component choose an element of
$\mathfrak q\setminus\mathfrak p$ and a power killing that
component. The product $a$ of these powers is outside
$\mathfrak p$ and satisfies $d_M(av)=0$.
The ideal $\operatorname{Ann}_A(v)$ is $\mathfrak p$-primary:
elements outside $\mathfrak p$ act invertibly, and a power of
$\mathfrak p$ kills $v$. Since $A$ is regular, hence
Cohen--Macaulay, this ideal contains an $A$-regular sequence of
length $h$. Explicitly, after choosing $i<h$ terms, the quotient
is Cohen--Macaulay and its associated primes have height $i$ in
$A$. The height-$h$ annihilator is contained in none of them; prime
avoidance supplies the next nonzerodivisor. This proves the assertion
inductively, using the Cohen--Macaulay quotient property
\cite[Lemma 10.104.2, tag 02JN]{StacksCMRegularSequence}.
Choose such a sequence and its ideal $J$.
The map $N=A/J\to Z^h(M)$ sending $1$ to $av$ extends by
(I8.1) to $C_J\to M$.

The $\mathfrak p$-component $z_{\mathfrak p}$ of the image of
$1\in N$ in $C_J^h$ generates
$C_{J,\mathfrak p}^h\simeq N_{\mathfrak p}$. The extension sends
$z_{\mathfrak p}$ to $av$, and therefore sends
$z_{\mathfrak p}/a$ to $v$. Thus every vector is in the image of
one such morphism. Taking the sum of these maps gives a surjection
of Cousin complexes
\[
                    P:=\bigoplus_\lambda C_{J_\lambda}
                    \longrightarrow M.                  \tag{I8.2}
\]

\emph{All residual endomorphisms act on the quotient.}
Every $C_J=K[J]$ is preserved by all $A$-linear graded
endomorphisms of $K$. Hence $P$ is a discrete continuous
$\End_A^*(K)$-module: a finite set of vectors of $P$ involves only
finitely many inputs in copies of $K$.
The kernel of (I8.2) is preserved by $B$ and by $\theta$ because
the map is a morphism of Cousin complexes. It is consequently
preserved by every word in these operators. It is also preserved by
every $f\in\End_A^*(K)$. To see this without assuming any
relation theorem, take one vector in the kernel, involving
finitely many inputs $z_1,\ldots,z_r$ in copies of $K$.
Proposition~\ref{i:dense} supplies a word polynomial $w$ in $B$
and $\theta$ agreeing with $f$ on all $z_i$. Thus the two operators
give exactly the same value on this vector, and $w$ preserves the
kernel. So $f$ preserves it as well.

It follows that $M$ has a continuous action of the whole
$\End_A^*(K)$ extending its diagonals and differential.
Continuity follows by lifting any finite set of vectors through
(I8.2) and using the corresponding finite set of inputs in $K$.
Uniqueness follows from density. The same density argument shows
that every morphism of Cousin complexes commutes with this action.
This construction proves all relations on every discrete Cousin
representation; it uses density to preserve a kernel in a known
discrete representation, not as a substitute for a presentation.

\emph{Passage to the complete target algebra.}
Let $Q$ be a target in the statement. Replace it by the closed
subalgebra generated by $B$ and $T$. Insert finite sums of the
point idempotents between factors in a word and pass to their
limits. The incidence relations show that a word of degree greater
than $\dim A$ vanishes. Thus this closed subalgebra has finitely
many nonzero graded pieces. Its topology still has a basis of open
left ideals. Replacing an open left ideal $L$ by
$\bigoplus_i(L\cap Q^i)$ gives a smaller open homogeneous left
ideal, since there are finitely many grades. These may therefore
be used as a basis.

For such an ideal $L$, the discrete graded left module $Q/L$
decomposes under $B$ into point-local torsion summands. In fact the
annihilator in $B$ of each vector is open, so it contains the full
factors outside a finite set of primes and suitable powers of the
maximal ideals at the remaining primes. The operator $T$ is its
differential. Decomposing by the constant offset
\emph{graded degree minus prime height} expresses $Q/L$ as a sum
of shifted Cousin complexes. By the result just proved it has a
unique continuous $\End_A^*(K)$-action extending its given
$B,T$-action.

Define
\[
 \phi_L:\End_A^*(K)\longrightarrow Q/L,
 \qquad f\longmapsto f(1+L).
\]
These maps are continuous and compatible as $L$ shrinks, because
the quotient maps are morphisms of the Cousin representations.
Completeness and separatedness give
$Q\simeq\varprojlim_L Q/L$ as topological additive groups.
The $\phi_L$ therefore define a continuous graded map
$\phi:\End_A^*(K)\to Q$, sending $B$ to its prescribed image
and $\theta$ to $T$. It is multiplicative: this holds on word
polynomials, which are dense, and both multiplications and $\phi$
are continuous. Density proves uniqueness as well.

Finally the finite topology on $\End_A^*(K)$ has a basis of open
homogeneous left ideals, namely annihilators of finite sets of
homogeneous inputs. It is separated and complete as proved above.
Thus it belongs to the target category and has exactly the asserted
universal property.
\end{proof}

\paragraph{Completion convention after I8.}
Theorem~\ref{i:linearuniversal} closes every higher-dimensional
relation and quotient-topology issue for the completion whose
neighborhoods are open left ideals, the intrinsic finite topology of
an endomorphism ring. This is stronger than an adelic action or dense
generation statement: it constructs the unique continuous map to
every complete target in that category. It does not identify an
unspecified completion with this one. The transcription of 43:66
says separated complete graded topological algebra but does not
unambiguously state the open-left-ideal convention; 43:64 also
mentions several possibly incompatible linear topologies. If
``topological algebra'' there is instead taken to allow arbitrary
topologies without a basis of open left ideals, that enlarged
universal problem is not settled by I8. The surface proof I2--I3
already works in that enlarged category.
